\section{Qué falta entre los Foundation Models y la Robotics}

Un foundation model sobresale en high-level reasoning, mientras que un sistema robótico debe producir control continuo en un plazo de decenas a cientos de milisegundos. Entre ambos no existe una API natural: difieren la representación de las entradas y salidas del modelo, la escala de los datos de entrenamiento, la frecuencia de control y las restricciones de seguridad. Esta sección desarrolla por completo el «por qué no se pueden conectar directamente» de la lecture.

\subsection{High-level reasoning y low-level control}

En la figura, PaLM-SayCan, a la izquierda, se encarga de elegir la skill de alto nivel a partir de un objetivo expresado en lenguaje; RT-1, a la derecha, se encarga de convertir la imagen y la instrucción en acciones de bajo nivel. El primero responde «qué hacer a continuación»; el segundo, «cómo deben moverse en este instante las articulaciones o el efector final». El plan de alto nivel tolera respuestas del orden de segundos y pasos discretos; la policy de bajo nivel debe manejar errores continuos y correcciones en lazo cerrado.

\lecturefigure{slide-09-foundation-models-robotics-gap.jpg}{Foundation Models + Robotics: high-level reasoning y low-level control}{Diapositiva V009 recuperada de la grabación oficial; 00:13:42}

Una policy de bajo nivel puede escribirse como
$$
a_t\sim \pi_\theta(a_t\mid o_{\le t},g),
$$
donde $a_t$ es la acción actual, $o_{\le t}$ es el observation history visual, de proprioception, etc., hasta el instante $t$, $g$ es el goal expresado en lenguaje y $\pi_\theta$ es la policy con parámetros $\theta$. A diferencia de una respuesta de texto que se produce una sola vez, esta distribución se vuelve a invocar en cada ciclo de control y modifica la observation del paso siguiente.

\begin{importantbox}{First use: policy y trajectory}
Una \term{Policy (política)} es la correspondencia entre la información actual y una distribución sobre acciones; una \term{trajectory (trayectoria)} es la secuencia temporal $(o_0,a_0,o_1,a_1,\ldots)$. La unidad básica de un robot dataset suele ser la trajectory, no un pair independiente de imagen y texto. Es más costosa, porque cada dato requiere que un agent real o simulado interactúe con el entorno en lazo cerrado.
\end{importantbox}

\subsection{Del modular stack al joint scaling}

La línea de tiempo pone en paralelo la percepción, planeación y control modulares (perception/planning/control) de 2014--2021 con RT-1, PaLM-E y RT-2 de 2021--2023. La tendencia no es que todos los módulos desaparezcan de golpe, sino que cada vez más representaciones y tareas se colocan en un espacio de parámetros compartido, lo que logra la model consolidation mediante mezclas de datos más grandes. La siguiente diapositiva resume la scaling recipe con «combine large diverse offline datasets», al tiempo que deja abierta la duda con «Scale up and done?».

\lecturefigure{slide-10-robotics-scaling-timeline.jpg}{Línea de tiempo de los sistemas robóticos: de lo modular a la model consolidation}{Diapositiva V010 recuperada de la grabación oficial; 00:15:06}

\lecturefigure{slide-11-scaling-recipe.jpg}{Scaling recipe: datos fuera de línea, modelo multitarea y positive transfer}{Diapositiva V011 recuperada de la grabación oficial; 00:15:36}

\begin{warningbox}{«Scale up» no es la respuesta predeterminada}
Los datos de lenguaje pueden recolectarse pasivamente de internet; los datos robóticos implican costos de embodiment, entorno, operator, reset y seguridad. Aumentar los parámetros solo puede producir transfer si la data mixture aporta supervision pertinente, si la action representation puede aprenderse y si la capacidad del modelo no está ocupada por tareas equivocadas. La escala es una variable experimental, no un lema que sustituya la definición del problema.
\end{warningbox}

\subsection{Moravec's paradox: la fluidez verbal no equivale a destreza manual}

Esta subsección explica con más detalle por qué el high-level model success no puede extrapolarse directamente al control. La Moravec's paradox señala que el symbolic reasoning, que a los humanos les parece de alto nivel, puede resultar fácil de implementar en una computadora, mientras que las capacidades sensoriomotoras, aunque subjetivamente «no requieren pensar», encierran una larga evolución y experiencia corporal. En la sujeción robótica, los errores milimétricos, los cambios de contacto, los objetos flexibles y las oclusiones suelen ser más difíciles que generar una explicación.

\lecturefigure{slide-12-moravec-paradox.jpg}{Moravec's paradox: la percepción y el movimiento de bajo nivel pueden ser más difíciles que el razonamiento de alto nivel}{Diapositiva V012 recuperada de la grabación oficial; 00:18:33}

\teachervoice{Fei Xia recurre a la Moravec's paradox para explicar por qué el éxito de los LLM no resolvió automáticamente la robótica: saber describir «cómo sujetar una taza» y poder sujetarla de forma estable ante variaciones de fricción, oclusiones e impactos de contacto son capacidades de niveles distintos. El curso toma esta brecha como el problema que define la low-level embodied intelligence.}

\subsection{Dos causas de fondo: escasez de datos e interfaz desalineada}

La primera challenge slide muestra, mediante diferencias de órdenes de magnitud, la escasez de robot trajectories; la segunda señala que un LLM produce token de texto, mientras que el robot necesita posición, rotación, gripper, velocidad o par. Aunque el LLM posea conocimiento semántico, si durante el entrenamiento no vio supervisión de acciones que pueda alinearse, o si su espacio de salida no puede expresar la precisión del control, no puede convertirse directamente en una policy.

\lecturefigure{slide-13-challenge-lack-of-data.jpg}{Challenge 1: hay muchos menos robot data que image-text / text data}{Diapositiva V013 recuperada de la grabación oficial; 00:20:21}

\lecturefigure{slide-14-challenge-llm-interface.jpg}{Challenge 2: el LLM no tiene una low-level control interface natural}{Diapositiva V014 recuperada de la grabación oficial; 00:21:30}

\lecturefigure{slide-15-agenda-two-parts.jpg}{Dos rutas: model consolidation y new interfaces}{Diapositiva V015 recuperada de la grabación oficial; 00:23:27}

\begin{knowledgebox}{Qué corrige cada una de las dos rutas}
La ruta de PaLM-E / RT-2 corrige principalmente el \textbf{reparto de datos y representaciones}: entrena conjuntamente las web-scale semantic tasks y las robot trajectories en un mismo modelo. La ruta de Language to Rewards corrige principalmente la \textbf{interfaz desalineada}: el modelo de lenguaje no produce directamente cada actuator command, sino una reward specification que un controlador puede optimizar.
\end{knowledgebox}

\subsection{Resumen de la sección}

La brecha entre los foundation models y el robot control puede reducirse a data regime, action representation, control loop y safety constraint. A continuación, la lecture intenta primero integrar el modelo y los datos en un VLA unificado y, después, conservar el controller tradicional y dejar al LLM solo a cargo de la reward interface.
